% GENERATED FILE. Edit canonical lessons, then run npm run generate:books.
% canonical-source-hash: fnv1a64:fea425cbada453b0
% canonical-lessons: JA-C163-kirei, JA-C163-shizuka, JA-C163-benri, JA-C163-rusu, JA-C163-seiketsu, JA-R163-review-158-163

\chapter{A Clean, Quiet and Handy Home}
\label{ch:ja-a-clean-quiet-and-handy-home}
\hlchaptermodality{163}

\begin{chapteropening}
\textbf{By the end of this chapter:} \emph{I can describe what my home is like.}

Complete the last lesson of chapter 163: I can describe what my home is like.
\end{chapteropening}

\section[\texorpdfstring{\ja{きれい} (kirei)}{きれい (kirei)}]{\ja{きれい} (kirei) — clean; pretty}

\label{lesson:JA-C163-kirei}



\begin{quote}
Before the new one: say the Japanese for repair, then the Japanese for old.
\end{quote}

\subsection*{You'll want to know: \ja{きれい}}
\textbf{\ja{きれい}} — \emph{kirei} — \textquotedblleft{}clean; pretty\textquotedblright{}.

\begin{quote}
\textbf{\ja{へやは} \ja{いつも} \ja{きれいです。}} — \emph{heya wa itsumo kirei desu} — \textquotedblleft{}The room is always clean.\textquotedblright{}
\end{quote}

\begin{cousinweb}
\textbf{\ja{きれい}} (\emph{kirei}) is Sino-Japanese (kango): two Chinese characters read with their on'yomi, readings borrowed from Chinese. They are \emph{ki}, \textquotedblleft{}figured cloth, beautiful\textquotedblright{} (Middle Chinese \emph{khjeX}) and \emph{rei}, \textquotedblleft{}lovely\textquotedblright{} (\emph{lejH}) (KANJIDIC2; Middle Chinese in Baxter's transcription). It is usually written in kana (JMdict).
\end{cousinweb}

\begin{grammarlens}[title={What you've built}]
One more word for everyday life at A2.
\end{grammarlens}

\subsection*{Your turn}
\begin{itemize}
\raggedright
  \item \emph{Say it:} \emph{kirei}
  \item \emph{Say it:} \emph{kirei}, once more
  \item \emph{Say it:} \emph{heya wa itsumo kirei desu} — The room is always clean.
  \item \emph{Recall:} say the Japanese for repair, then the Japanese for old, then say \emph{kirei} again
\end{itemize}

\begin{tcolorbox}[breakable,colback=teal!4,colframe=teal!35!black,title={Before you move on}]
What does \ja{きれい} mean? (\textquotedblleft{}Clean; pretty\textquotedblright{}.) Say it once more.
\end{tcolorbox}
\section[\texorpdfstring{\ja{しずか} (shizuka)}{しずか (shizuka)}]{\ja{しずか} (shizuka) — quiet; calm}

\label{lesson:JA-C163-shizuka}



\begin{quote}
Before the new one: say the Japanese for old, then the Japanese for clean.
\end{quote}

\subsection*{You'll want to know: \ja{しずか}}
\textbf{\ja{しずか}} — \emph{shizuka} — \textquotedblleft{}quiet; calm\textquotedblright{}.

\begin{quote}
\textbf{\ja{この} \ja{まちは} \ja{しずかです。}} — \emph{kono machi wa shizuka desu} — \textquotedblleft{}This town is quiet.\textquotedblright{}
\end{quote}

\begin{cousinweb}
\textbf{\ja{しずか}} (\emph{shizuka}) is a native Japanese word (wago), inherited rather than borrowed from Chinese; its kanji was attached later as a kun'yomi, and its older origin is not settled (JMdict, KANJIDIC2).
\end{cousinweb}

\begin{grammarlens}[title={What you've built}]
One more word for everyday life at A2.
\end{grammarlens}

\subsection*{Your turn}
\begin{itemize}
\raggedright
  \item \emph{Say it:} \emph{shizuka}
  \item \emph{Say it:} \emph{shizuka}, once more
  \item \emph{Say it:} \emph{kono machi wa shizuka desu} — This town is quiet.
  \item \emph{Recall:} say the Japanese for old, then the Japanese for clean, then say \emph{shizuka} again
\end{itemize}

\begin{tcolorbox}[breakable,colback=teal!4,colframe=teal!35!black,title={Before you move on}]
What does \ja{しずか} mean? (\textquotedblleft{}Quiet; calm\textquotedblright{}.) Say it once more.
\end{tcolorbox}
\section[\texorpdfstring{\ja{べんり} (benri)}{べんり (benri)}]{\ja{べんり} (benri) — convenient; handy}

\label{lesson:JA-C163-benri}



\begin{quote}
Before the new one: say the Japanese for clean, then the Japanese for quiet.
\end{quote}

\subsection*{You'll want to know: \ja{べんり}}
\textbf{\ja{べんり}} — \emph{benri} — \textquotedblleft{}convenient; handy\textquotedblright{}.

\begin{quote}
\textbf{\ja{この} \ja{だいどころは} \ja{べんりです。}} — \emph{kono daidokoro wa benri desu} — \textquotedblleft{}This kitchen is convenient.\textquotedblright{}
\end{quote}

\begin{cousinweb}
\textbf{\ja{べんり}} (\emph{benri}) is Sino-Japanese (kango): two Chinese characters read with their on'yomi, readings borrowed from Chinese. They are \emph{ben}, \textquotedblleft{}convenience\textquotedblright{} (Middle Chinese \emph{bjienH}) and \emph{ri}, \textquotedblleft{}profit, advantage\textquotedblright{} (\emph{lijH}) (KANJIDIC2; Middle Chinese in Baxter's transcription).
\end{cousinweb}

\begin{grammarlens}[title={What you've built}]
One more word for everyday life at A2.
\end{grammarlens}

\subsection*{Your turn}
\begin{itemize}
\raggedright
  \item \emph{Say it:} \emph{benri}
  \item \emph{Say it:} \emph{benri}, once more
  \item \emph{Say it:} \emph{kono daidokoro wa benri desu} — This kitchen is convenient.
  \item \emph{Recall:} say the Japanese for clean, then the Japanese for quiet, then say \emph{benri} again
\end{itemize}

\begin{tcolorbox}[breakable,colback=teal!4,colframe=teal!35!black,title={Before you move on}]
What does \ja{べんり} mean? (\textquotedblleft{}Convenient; handy\textquotedblright{}.) Say it once more.
\end{tcolorbox}
\section[\texorpdfstring{\ja{るす} (rusu)}{るす (rusu)}]{\ja{るす} (rusu) — being out; being away from home}

\label{lesson:JA-C163-rusu}



\begin{quote}
Before the new one: say the Japanese for quiet, then the Japanese for convenient.
\end{quote}

\subsection*{You'll want to know: \ja{るす}}
\textbf{\ja{るす}} — \emph{rusu} — \textquotedblleft{}being out; being away from home\textquotedblright{}.

\begin{quote}
\textbf{\ja{どようびは} \ja{いつも} \ja{るすです。}} — \emph{doyōbi wa itsumo rusu desu} — \textquotedblleft{}On Saturdays I am always out.\textquotedblright{}
\end{quote}

\begin{cousinweb}
\textbf{\ja{るす}} (\emph{rusu}) is Sino-Japanese (kango): two Chinese characters read with their on'yomi, readings borrowed from Chinese. They are \emph{ru}, \textquotedblleft{}detain, stop\textquotedblright{} (Middle Chinese \emph{ljuw}) and \emph{su}, \textquotedblleft{}guard\textquotedblright{} (\emph{syuwX}) (KANJIDIC2; Middle Chinese in Baxter's transcription).
\end{cousinweb}

\begin{grammarlens}[title={What you've built}]
One more word for everyday life at A2.
\end{grammarlens}

\subsection*{Your turn}
\begin{itemize}
\raggedright
  \item \emph{Say it:} \emph{rusu}
  \item \emph{Say it:} \emph{rusu}, once more
  \item \emph{Say it:} \emph{doyōbi wa itsumo rusu desu} — On Saturdays I am always out.
  \item \emph{Recall:} say the Japanese for quiet, then the Japanese for convenient, then say \emph{rusu} again
\end{itemize}

\begin{tcolorbox}[breakable,colback=teal!4,colframe=teal!35!black,title={Before you move on}]
What does \ja{るす} mean? (\textquotedblleft{}Being out; being away from home\textquotedblright{}.) Say it once more.
\end{tcolorbox}
\section[\texorpdfstring{\ja{せいけつ} (seiketsu)}{せいけつ (seiketsu)}]{\ja{せいけつ} (seiketsu) — hygienic; clean and tidy}

\label{lesson:JA-C163-seiketsu}



\begin{quote}
Before the new one: say the Japanese for convenient, then the Japanese for being out.
\end{quote}

\subsection*{You'll want to know: \ja{せいけつ}}
\textbf{\ja{せいけつ}} — \emph{seiketsu} — \textquotedblleft{}hygienic; clean and tidy\textquotedblright{}.

\begin{quote}
\textbf{\ja{だいどころは} \ja{いつも} \ja{せいけつです。}} — \emph{daidokoro wa itsumo seiketsu desu} — \textquotedblleft{}The kitchen is always hygienic.\textquotedblright{}
\end{quote}

\begin{cousinweb}
\textbf{\ja{せいけつ}} (\emph{seiketsu}) is Sino-Japanese (kango): two Chinese characters read with their on'yomi, readings borrowed from Chinese. They are \emph{sei}, \textquotedblleft{}pure, clear\textquotedblright{} (Middle Chinese \emph{tshjeng}) and \emph{ketsu}, \textquotedblleft{}clean, pure\textquotedblright{} (\emph{ket}) (KANJIDIC2; Middle Chinese in Baxter's transcription).
\end{cousinweb}

\begin{grammarlens}[title={What you've built}]
One more word for everyday life at A2.
\end{grammarlens}

\subsection*{Your turn}
\begin{itemize}
\raggedright
  \item \emph{Say it:} \emph{seiketsu}
  \item \emph{Say it:} \emph{seiketsu}, once more
  \item \emph{Say it:} \emph{daidokoro wa itsumo seiketsu desu} — The kitchen is always hygienic.
  \item \emph{Recall:} say the Japanese for convenient, then the Japanese for being out, then say \emph{seiketsu} again
\end{itemize}

\begin{tcolorbox}[breakable,colback=teal!4,colframe=teal!35!black,title={Before you move on}]
What does \ja{せいけつ} mean? (\textquotedblleft{}Hygienic; clean and tidy\textquotedblright{}.) Say it once more.
\end{tcolorbox}
\section[(dialogue)]{Words from chapters 158–163 — a review}

\label{lesson:JA-R163-review-158-163}



\begin{quote}
Say the Japanese for a key and to switch off.
\end{quote}

\subsection*{Your turn}
Say these aloud:

\begin{itemize}
\raggedright
  \item a key, to switch off, electric light, to open up, dark
  \item housework, to throw away, to sort, to stack, sorting out
  \item cleaning, dust specks, to get dirty, to make dirty, dirty
  \item washing clothes, to get wet, to spread out, to get dry, to dry
  \item to be broken, to shatter, to come off, repair, old
  \item clean, quiet, convenient, being out, hygienic
\end{itemize}

\begin{tcolorbox}[breakable,colback=teal!4,colframe=teal!35!black,title={Before you move on}]
A key? (\textbf{\ja{かぎ}}, \emph{kagi}.) To switch off? (\textbf{\ja{けす}}, \emph{kesu}.) Electric light? (\textbf{\ja{でんき}}, \emph{denki}.) To open up? (\textbf{\ja{あける}}, \emph{akeru}.) Dark? (\textbf{\ja{くらい}}, \emph{kurai}.) Housework? (\textbf{\ja{かじ}}, \emph{kaji}.) To throw away? (\textbf{\ja{すてる}}, \emph{suteru}.) To sort? (\textbf{\ja{わける}}, \emph{wakeru}.) To stack? (\textbf{\ja{かさねる}}, \emph{kasaneru}.) Sorting out? (\textbf{\ja{せいり}}, \emph{seiri}.) Cleaning? (\textbf{\ja{そうじ}}, \emph{sōji}.) Dust specks? (\textbf{\ja{ちり}}, \emph{chiri}.) To get dirty? (\textbf{\ja{よごれる}}, \emph{yogoreru}.) To make dirty? (\textbf{\ja{よごす}}, \emph{yogosu}.) Dirty? (\textbf{\ja{きたない}}, \emph{kitanai}.) Washing clothes? (\textbf{\ja{せんたく}}, \emph{sentaku}.) To get wet? (\textbf{\ja{ぬれる}}, \emph{nureru}.) To spread out? (\textbf{\ja{ひろげる}}, \emph{hirogeru}.) To get dry? (\textbf{\ja{かわく}}, \emph{kawaku}.) To dry? (\textbf{\ja{かわかす}}, \emph{kawakasu}.) To be broken? (\textbf{\ja{こわれる}}, \emph{kowareru}.) To shatter? (\textbf{\ja{われる}}, \emph{wareru}.) To come off? (\textbf{\ja{とれる}}, \emph{toreru}.) Repair? (\textbf{\ja{しゅうり}}, \emph{shūri}.) Old? (\textbf{\ja{ふるい}}, \emph{furui}.) Clean? (\textbf{\ja{きれい}}, \emph{kirei}.) Quiet? (\textbf{\ja{しずか}}, \emph{shizuka}.) Convenient? (\textbf{\ja{べんり}}, \emph{benri}.) Being out? (\textbf{\ja{るす}}, \emph{rusu}.) Hygienic? (\textbf{\ja{せいけつ}}, \emph{seiketsu}.)
\end{tcolorbox}
